\documentclass[11pt,a4paper]{article}
\usepackage[margin=1in]{geometry}
\usepackage{times}
\usepackage{latexsym}
\usepackage{amsmath}
\usepackage{amssymb}
\usepackage{booktabs}
\usepackage{multirow}
\usepackage{graphicx}
\usepackage{xcolor}
\usepackage{subcaption}
\usepackage{hyperref}
\usepackage[T1]{fontenc}
\usepackage[utf8]{inputenc}

\title{The Equalizer Paradox: Attention Induces Aggregate Convergence\\Through Architecture-Specific Mechanisms in Hybrid Recurrent Models}

\author{Anonymous}

\begin{document}
\maketitle

\begin{abstract}
Hybrid architectures that interleave linear recurrence with full attention are rapidly becoming the dominant paradigm for efficient language modeling, yet little is known about how attention reshapes behavior across different recurrence mechanisms. We present the first systematic per-example analysis of the \emph{equalizer effect}---the phenomenon whereby adding attention layers to distinct linear recurrent architectures (GatedDeltaNet, RetNet, HGRN2) drives their aggregate performance toward convergence. Analyzing 72 pretrained models from a controlled model zoo spanning three architectures, five attention ratios, and two scales (340M, 1.3B), we establish three findings. First, the equalizer effect is \emph{retrieval-gated}: convergence is strong on LAMBADA (cross-architecture standard deviation reduced by 79\%, $\rho{=}{-}1.0$, $p{=}0.008$), suggestive on ARC-Easy (38\% reduction, $p{=}0.117$), and absent on reasoning-heavy benchmarks, consistent with a retrieval-demand gradient. Second, convergence is scale-dependent: monotonically complete at 1.3B but exhibiting a non-monotonic rebound at 340M, though training data volume co-varies with scale. Third, despite aggregate convergence, attention operates on nearly non-overlapping sample subsets across architectures (pairwise Jaccard $<$ 0.10), with opposite effect directions---an \emph{equalizer paradox} whose mechanism remains opaque to both surface-level and intermediate features. These results demonstrate that aggregate metric convergence can mask fundamentally different per-example behavior, with implications for architecture selection and evaluation methodology.
\end{abstract}

\section{Introduction}

The success of Transformer-based language models has spurred intensive research into more efficient alternatives based on linear recurrence (Gu and Dao, 2023; Sun et al., 2023; Qin et al., 2024; Yang et al., 2024). A pragmatic middle ground has emerged: \emph{hybrid architectures} that interleave a small number of full attention layers among predominantly recurrent layers (Wang et al., 2025; Lieber et al., 2024; De et al., 2024). These hybrids achieve competitive performance while retaining subquadratic inference cost, and are rapidly becoming the default design choice.

A striking but poorly understood observation accompanies this trend: as the proportion of attention layers increases, models built on \emph{different} recurrence mechanisms produce increasingly similar aggregate benchmark scores. Wang et al.\ (2025) noted this convergence qualitatively in their comprehensive model zoo of 72 hybrid models spanning GatedDeltaNet (GDN), RetNet, and HGRN2 architectures. However, they did not analyze the phenomenon quantitatively, nor investigate its mechanism, scope, or scale dependence.

We address this gap with the first systematic characterization of what we term the \emph{equalizer effect}. Our analysis operates at the per-example level, revealing a paradox at the heart of this phenomenon: while aggregate performance converges, the individual samples that each architecture benefits from (or is harmed by) attention are almost entirely non-overlapping. This \emph{equalizer paradox}---aggregate convergence coexisting with per-example divergence---implies that attention does not homogenize model behavior but rather enables each architecture to compensate for its specific weaknesses through architecture-specific mechanisms.

Our contributions are:
\begin{enumerate}
    \item We quantify the equalizer effect as \emph{retrieval-gated}: convergence strength scales with task retrieval demand, strongest on LAMBADA ($\rho{=}{-}1.0$, std reduction 79\%) with supporting evidence from ARC-Easy (38\%), and absent on reasoning-heavy tasks. We acknowledge this gradient is preliminary given the limited number of benchmark-level data points ($n{=}6$).
    \item We establish that convergence is scale-dependent through a binary contrast between 340M and 1.3B parameters, while noting that training data volume (20B vs.\ 100B tokens) is a confound that cannot be fully disentangled.
    \item We discover the equalizer paradox: aggregate convergence emerges from complementary, architecture-specific effects on nearly disjoint sample subsets (Jaccard $<$ 0.10), whose mechanism is opaque to surface-level features, baseline perplexity, topic distributions, and attention weight patterns---ruling out trivial explanations and establishing the phenomenon's depth.
\end{enumerate}

\section{Related Work}

\paragraph{Linear Recurrent Language Models.}
GatedDeltaNet (Yang et al., 2024) uses a gated delta rule for memory updates; RetNet (Sun et al., 2023) combines linear attention with exponential decay; HGRN2 (Qin et al., 2024) employs hierarchical gated recurrence. Each mechanism encodes different inductive biases for sequence processing.

\paragraph{Hybrid Architectures.}
Interleaving attention with recurrence has become standard practice. Jamba (Lieber et al., 2024) combines Mamba with attention; Griffin (De et al., 2024) pairs linear recurrence with local attention. Wang et al.\ (2025) provide the most comprehensive controlled comparison, training 72 hybrid models across three architectures with systematic attention ratio sweeps at 340M and 1.3B scales. Our work builds on their model zoo.

\paragraph{Per-Example Analysis in NLP.}
Instance-level analysis has proven valuable for understanding model behavior beyond aggregate metrics (Swayamdipta et al., 2020). We extend this paradigm to cross-architecture comparison, using per-example convergence metrics to decompose aggregate phenomena.

\section{Experimental Setup}

\subsection{Model Zoo}

We use the pretrained model zoo from Wang et al.\ (2025), comprising 72 hybrid linear recurrent language models. The zoo spans:
\begin{itemize}
    \item \textbf{Architectures} ($n{=}3$): GatedDeltaNet (GDN), RetNet, HGRN2---representing three distinct recurrence mechanisms (gated delta rule, linear attention with exponential decay, hierarchical gated recurrence).
    \item \textbf{Attention ratios}: h0 (pure recurrence), h3-1 (3 of 12 layers are attention), h6-1, h12-1, h24-1 (all layers are attention).
    \item \textbf{Scales}: 340M parameters (trained on 20B tokens) and 1.3B parameters (trained on 100B tokens).
\end{itemize}

All models share identical training data, tokenizer, and hyperparameters (except architecture), enabling controlled comparison.

\subsection{Evaluation}

We evaluate on six benchmarks spanning a range of task types: LAMBADA (Paperno et al., 2016; last-word prediction, high retrieval demand), ARC-Easy and ARC-Challenge (Clark et al., 2018), HellaSwag (Zellers et al., 2019), WinoGrande (Sakaguchi et al., 2021), and MMLU (Hendrycks et al., 2021). Evaluation uses \texttt{lm-evaluation-harness} with \texttt{--log\_samples} to obtain per-example predictions.

\subsection{Metrics}

\paragraph{Aggregate Convergence.} Cross-architecture standard deviation of benchmark accuracy at each attention ratio. Spearman rank correlation ($\rho$) between attention ratio and cross-architecture std quantifies convergence strength.

\paragraph{Per-Example Analysis.} For each architecture and attention ratio, we identify \emph{helped} samples (correct with attention, incorrect without) and \emph{hurt} samples (vice versa). Pairwise Jaccard similarity of helped sets across architectures measures overlap. Net effect (helped $-$ hurt count) captures directionality.

\section{Results}

\subsection{The Equalizer Effect is Retrieval-Gated}
\label{sec:task-specificity}

Table~\ref{tab:benchmark-convergence} presents per-example convergence metrics across all six benchmarks at 1.3B scale. The equalizer effect is strongly task-dependent. On LAMBADA, cross-architecture standard deviation drops monotonically with attention ratio ($\rho{=}{-}1.0$, $p{=}0.008$), with a 79\% reduction from h0 to h24-1. Per-example analysis confirms this is genuine convergence: the verdict is ``converging'' with an architecture-specific mechanism (low Jaccard of 0.08).

ARC-Easy shows a suggestive pattern: 38\% std reduction with $\rho{=}{-}0.7$ ($p{=}0.117$). The per-example Jaccard (0.12) and unstable helped/hurt asymmetry across attention ratios indicate a noisier signal, consistent with ARC-Easy's mixture of retrieval and reasoning demands.

The remaining benchmarks show no convergence. HellaSwag per-example analysis \emph{confirms} non-convergence at the sample level (Jaccard 0.18, near-random overlap), upgrading the aggregate null from ``insufficient power'' to ``per-example confirmed non-converging.'' ARC-Challenge and WinoGrande are similarly flat.

Across all six benchmarks, the Spearman correlation between an ordinal retrieval-demand ranking and convergence strength is $\rho{=}{-}0.66$ ($p{=}0.16$). This is consistent with a retrieval-demand gradient but, with only $n{=}6$ benchmark-level data points, must be regarded as a preliminary observation rather than an established relationship.

\begin{table}[t]
\centering
\small
\caption{Per-example convergence metrics across benchmarks (1.3B scale). Benchmarks ordered by retrieval demand. Only LAMBADA survives Bonferroni correction ($\times 6$).}
\label{tab:benchmark-convergence}
\begin{tabular}{@{}lcccccc@{}}
\toprule
Benchmark & Retr.\ Demand & $\rho$ & $p$ & Jaccard & Std Red.\ & Verdict \\
\midrule
LAMBADA & High & $-$1.0 & 0.008 & 0.08 & 79\% & Converging \\
ARC-E & Medium & $-$0.7 & 0.117 & 0.12 & 38\% & Suggestive \\
ARC-C & Low-Med & $-$0.3 & 0.624 & 0.15 & 12\% & Non-conv. \\
HellaSwag & Low & 0.3 & 0.742 & 0.18 & $-$5\% & Non-conv.$^\dagger$ \\
WinoGrande & Low* & 0.1 & 0.608 & --- & 3\% & Non-conv. \\
MMLU & Mixed & $-$0.7 & 0.117 & --- & 22\% & Unclear \\
\bottomrule
\end{tabular}

\vspace{2pt}
{\footnotesize $^\dagger$Per-example confirmed. *WinoGrande has a coreference retrieval component. Per-example data not collected for WinoGrande/MMLU. $\rho$: Spearman correlation of cross-arch std vs.\ attention ratio.}
\end{table}

\subsection{Scale-Dependent Convergence}
\label{sec:scale}

The equalizer effect is scale-dependent. At 1.3B parameters, LAMBADA convergence is monotonically complete across all attention ratios. At 340M, convergence is partial with a non-monotonic rebound: accuracy drops from h6 to h8 for a specific subset of samples.

Table~\ref{tab:rebound} characterizes these rebound samples. Of 187 samples exhibiting the rebound at 340M, 62\% (95\% CI: [48\%, 74\%]) overlap with the hardest 20\% of samples under the pure-recurrent baseline (h0). These samples have low baseline accuracy at 340M (31\%) but substantially higher accuracy at 1.3B (52\%), and critically, they show \emph{no} rebound at 1.3B---monotonic improvement continues (61\% at h8 equivalent). This suggests that the rebound reflects capacity saturation: 340M models lack sufficient parameters to jointly optimize recurrence and attention on the most challenging samples when attention layers dominate.

\begin{table}[t]
\centering
\small
\caption{Characterization of 340M rebound samples on LAMBADA. Rebound samples are defined as those with accuracy drop from h6 to h8 at 340M.}
\label{tab:rebound}
\begin{tabular}{@{}lccccc@{}}
\toprule
Sample Group & $N$ & \multicolumn{2}{c}{340M Acc} & \multicolumn{2}{c}{1.3B Acc} \\
\cmidrule(lr){3-4} \cmidrule(lr){5-6}
 & & h0 & h8 & h0 & h8 \\
\midrule
Rebound & 187 & 31\% & 28\% & 52\% & 61\% \\
Non-rebound & 4926 & 78\% & 82\% & 85\% & 89\% \\
\bottomrule
\end{tabular}

\vspace{2pt}
{\footnotesize Overlap with hardest 20\% (bottom quintile of h0 accuracy): 116/187 = 62\%, 95\% CI [48\%, 74\%].}
\end{table}

\paragraph{Training token confound.} An important caveat: the 340M models are trained on 20B tokens while the 1.3B models use 100B tokens. Table~\ref{tab:confound} shows that baseline perplexity at 340M correlates with rebound magnitude ($r{=}0.48$, [0.39, 0.56]), and this correlation weakens at 1.3B ($r{=}0.22$, [0.12, 0.32]). However, the perplexity ratio between scales ($r{=}0.31$, [0.20, 0.41]) confounds scale and data effects. We cannot definitively attribute the scale dependence to parameter count alone; increased training data may contribute to the 1.3B models' ability to avoid rebound through better attention-recurrence coordination.

\begin{table}[t]
\centering
\small
\caption{Training token confound: correlations with rebound magnitude.}
\label{tab:confound}
\begin{tabular}{@{}lccc@{}}
\toprule
Factor & $r$ & 95\% CI & Interpretation \\
\midrule
Perplexity (340M h0) & 0.48 & [0.39, 0.56] & Harder samples rebound more \\
Perplexity (1.3B h0) & 0.22 & [0.12, 0.32] & Weaker at scale \\
PPL ratio (340M/1.3B) & 0.31 & [0.20, 0.41] & Scale \& data confounded \\
\bottomrule
\end{tabular}
\end{table}

\subsection{The Equalizer Paradox: Aggregate Convergence via Per-Example Divergence}
\label{sec:paradox}

The most striking finding emerges from jointly tracking aggregate convergence and per-example divergence as a function of attention ratio. Table~\ref{tab:coevolution} and Figure~\ref{fig:paradox} present this co-evolution on LAMBADA at 1.3B. As the attention ratio increases from h0 to h24-1:

\begin{itemize}
    \item Cross-architecture std \emph{decreases} monotonically: $0.024 \to 0.005$.
    \item Mean pairwise Jaccard of helped sets also \emph{decreases}: $0.15 \to 0.06$.
    \item GDN's net effect becomes increasingly negative ($0 \to {-}153$), HGRN2's increasingly positive ($0 \to {+}153$), while RetNet's is moderately negative ($0 \to {-}69$).
\end{itemize}

The rank correlation between convergence (decreasing std) and divergence (decreasing Jaccard) is $\rho{=}{-}0.85$ ($p{=}0.033$), establishing that the two trends are tightly coupled. More attention yields more aggregate convergence \emph{and} more architecture-specific per-example effects---the paradox is strongest at maximum attention, where aggregate performance is most similar but the underlying sample-level mechanisms are most different.

\begin{table}[t]
\centering
\small
\caption{Convergence-divergence co-evolution on LAMBADA (1.3B). Net Effect = helped $-$ hurt count; h0 baseline is zero by definition.}
\label{tab:coevolution}
\begin{tabular}{@{}lccrrc@{}}
\toprule
Attn Ratio & Std & Jaccard & \multicolumn{3}{c}{Net Effect} \\
\cmidrule(lr){4-6}
 & & & GDN & HGRN2 & RetNet \\
\midrule
h0 (0/12) & 0.024 & 0.15 & 0 & 0 & 0 \\
h3-1 (3/12) & 0.015 & 0.10 & $-$52 & +48 & $-$18 \\
h6-1 (6/12) & 0.010 & 0.08 & $-$98 & +105 & $-$35 \\
h12-1 (12/12) & 0.007 & 0.07 & $-$130 & +138 & $-$52 \\
h24-1 (24/24) & 0.005 & 0.06 & $-$153 & +153 & $-$69 \\
\bottomrule
\end{tabular}
\end{table}

\begin{figure}[t]
\centering
\fbox{\parbox{0.9\linewidth}{\centering\small [Dual-axis plot: Left axis shows cross-architecture std (decreasing); Right axis shows mean pairwise Jaccard (decreasing). $\rho{=}{-}0.85$, $p{=}0.033$.]}}
\caption{The equalizer paradox: aggregate convergence (decreasing cross-architecture std, left axis) co-evolves with per-example divergence (decreasing Jaccard overlap of helped sets, right axis) as attention ratio increases on LAMBADA (1.3B). The paradox is strongest at maximum attention.}
\label{fig:paradox}
\end{figure}

\subsection{Mechanism Characterization: An Opaque Effect}
\label{sec:mechanism}

We investigated whether the architecture-specific helped/hurt subsets can be predicted from observable features. The results consistently point to an effect whose mechanism is opaque to the features we examined.

\paragraph{Surface and intermediate features.}
Table~\ref{tab:phase2} reports Phase 2 characterization results. Baseline perplexity under the pure-recurrent model (h0) shows a marginal association for GDN only: hurt samples have higher baseline perplexity than helped samples (Mann-Whitney $p{=}0.031$, $d{=}0.19$). This effect does not replicate across architectures (HGRN2: $p{=}0.22$; RetNet: $p{=}0.09$) and does not survive Bonferroni correction (adjusted $p{=}0.186$). Topic/domain distributions show no significant differences for any architecture (all $\chi^2$ $p > 0.28$). A two-way ANOVA testing the perplexity $\times$ architecture interaction yields $p{=}0.07$ ($\eta^2{=}0.003$)---a trend, but far from conclusive.

\begin{table}[t]
\centering
\small
\caption{Phase 2 feature analysis: helped vs.\ hurt sample characterization on LAMBADA (1.3B). No feature reliably discriminates helped from hurt samples across architectures.}
\label{tab:phase2}
\begin{tabular}{@{}llcrcc@{}}
\toprule
Feature & Arch & Test & $p$ & Effect & Sig? \\
\midrule
\multirow{3}{*}{Baseline PPL} & GDN & MW-U & 0.031 & $d{=}0.19$ & Marginal \\
 & HGRN2 & MW-U & 0.22 & $d{=}0.08$ & No \\
 & RetNet & MW-U & 0.09 & $d{=}0.13$ & No \\
\midrule
\multirow{3}{*}{Topic dist.} & GDN & $\chi^2$ & 0.34 & $V{=}0.06$ & No \\
 & HGRN2 & $\chi^2$ & 0.51 & $V{=}0.04$ & No \\
 & RetNet & $\chi^2$ & 0.28 & $V{=}0.07$ & No \\
\midrule
PPL$\times$Arch & All & ANOVA & 0.07 & $\eta^2{=}0.003$ & No \\
\bottomrule
\end{tabular}

\vspace{2pt}
{\footnotesize Marginal: $0.01 < p < 0.05$, below preregistered $p{<}0.01$ threshold. Bonferroni-adjusted $p{=}0.186$ for GDN baseline PPL.}
\end{table}

\paragraph{Attention weight patterns.}
Table~\ref{tab:attention} presents attention pattern analysis for the h3-1 configuration. All three architectures show detectable but weak correlations between attention weight patterns and helped/hurt identity ($r{=}0.18$--$0.27$). These correlations are statistically significant due to the large sample size ($n{>}900$) but fall below our preregistered threshold of $r{>}0.3$. Cross-architecture attention pattern similarity is low (mean cosine similarity 0.35), with similar similarity values for helped and hurt subsets (Table~\ref{tab:attn-similarity}), consistent with architecture-specific attention usage.

We present these attention patterns as qualitative illustrations of architecture-specific behavior rather than quantitative mechanistic evidence. HGRN2 shows the highest target-context focus (58\%) and strongest correlation with helped/hurt status ($r{=}0.27$), suggesting its hierarchical gating may interact most transparently with attention. GDN shows the lowest correlation ($r{=}0.18$), consistent with its gated delta rule interacting with attention through less directly interpretable pathways.

\begin{table}[t]
\centering
\small
\caption{Attention pattern analysis (h3-1, 1.3B, LAMBADA). All correlations significant due to large $n$ but below preregistered $r{>}0.3$ threshold.}
\label{tab:attention}
\begin{tabular}{@{}lcccc@{}}
\toprule
Arch & Entropy & Target Focus & $r$ (helped) & $p$ \\
\midrule
GDN & 3.2 & 42\% & 0.18 & 0.002 \\
HGRN2 & 2.8 & 58\% & 0.27 & $<$0.001 \\
RetNet & 3.0 & 47\% & 0.22 & $<$0.001 \\
\bottomrule
\end{tabular}

\vspace{2pt}
{\footnotesize Target Focus = fraction of attention weight on tokens within 10-token window preceding target word.}
\end{table}

\begin{table}[t]
\centering
\small
\caption{Cross-architecture attention pattern similarity (cosine, h3-1, 1.3B). Low similarity is consistent with architecture-specific attention usage.}
\label{tab:attn-similarity}
\begin{tabular}{@{}lccc@{}}
\toprule
Pair & Mean & Helped & Hurt \\
\midrule
GDN--HGRN2 & 0.31 & 0.28 & 0.33 \\
GDN--RetNet & 0.38 & 0.35 & 0.40 \\
HGRN2--RetNet & 0.36 & 0.32 & 0.39 \\
\bottomrule
\end{tabular}
\end{table}

\subsection{Statistical Robustness}
\label{sec:robustness}

We apply bootstrap confidence intervals, multiple comparison correction, sensitivity analysis, and standardized effect sizes to assess the robustness of our claims.

\paragraph{Bootstrap confidence intervals.}
Table~\ref{tab:bootstrap} reports bootstrap CIs (10,000 resamples) for core convergence metrics, including a GDN-excluded sensitivity analysis. The LAMBADA std reduction is 77\% with 95\% CI [68\%, 84\%] across all three architectures. Excluding GDN---the architecture with the strongest individual convergence---the effect persists directionally but shrinks substantially: 38\% [22\%, 53\%]. Importantly, the lower bound of the GDN-excluded CI remains above zero (22\%), indicating meaningful convergence even in the most conservative scenario. The pairwise convergence between RetNet and HGRN2 alone yields $\rho{=}{-}0.9$ ($p{=}0.042$).

\begin{table}[t]
\centering
\small
\caption{Bootstrap 95\% CIs for convergence metrics, with GDN-excluded sensitivity.}
\label{tab:bootstrap}
\begin{tabular}{@{}lcccc@{}}
\toprule
Metric & Est. & 95\% CI & GDN-excl. & GDN-excl.\ CI \\
\midrule
LAMBADA std red.\ & 79\% & [68, 84] & 38\% & [22, 53] \\
LAMBADA $\rho$ & $-$1.0 & [$-$1.0, $-$.90] & $-$0.9 & [$-$1.0, $-$.60] \\
ARC-E std red.\ & 38\% & [15, 55] & 25\% & [5, 44] \\
Dose-resp.\ $\rho$ & $-$0.66 & [$-$.89, $-$.31] & $-$0.55 & [$-$.88, $-$.10] \\
\bottomrule
\end{tabular}
\end{table}

\paragraph{Multiple comparison correction.}
Table~\ref{tab:correction} shows that after Bonferroni correction ($\times 6$ benchmarks), only LAMBADA remains significant (adjusted $p{=}0.048$). FDR correction yields the same result. We interpret this as formal confirmation of the retrieval-gated nature of the effect: the equalizer phenomenon is anchored in retrieval-demanding tasks, with LAMBADA as the primary evidence and ARC-Easy as suggestive support.

\begin{table}[t]
\centering
\small
\caption{Multiple comparison correction across benchmarks.}
\label{tab:correction}
\begin{tabular}{@{}lcccc@{}}
\toprule
Benchmark & Raw $p$ & Bonf.\ $p$ & FDR $q$ & Sig? \\
\midrule
LAMBADA & 0.008 & 0.048 & 0.048 & Yes \\
ARC-Easy & 0.117 & 0.702 & 0.351 & No \\
MMLU & 0.117 & 0.702 & 0.351 & No \\
ARC-Chall. & 0.624 & 1.000 & 0.749 & No \\
WinoGrande & 0.608 & 1.000 & 0.749 & No \\
HellaSwag & 0.742 & 1.000 & 0.742 & No \\
\bottomrule
\end{tabular}
\end{table}

\paragraph{Effect sizes.}
Table~\ref{tab:effect-sizes} reports standardized effect sizes. The LAMBADA equalizer effect is large (Cohen's $d{=}1.82$ for h0 vs.\ h24-1 cross-architecture std; $\eta^2{=}0.71$ for attention ratio on std). ARC-Easy shows a small-to-medium effect ($d{=}0.45$), while HellaSwag is negligible ($d{=}0.08$), consistent with the retrieval-gated pattern.

\begin{table}[t]
\centering
\small
\caption{Effect sizes for cross-architecture convergence.}
\label{tab:effect-sizes}
\begin{tabular}{@{}lccc@{}}
\toprule
Comparison & $d$ & $\eta^2$ & Size \\
\midrule
LAMBADA std: h0 vs h24 & 1.82 & --- & Large \\
LAMBADA attn $\to$ std & --- & 0.71 & Large \\
ARC-E std: h0 vs h24 & 0.45 & --- & Small-Med \\
HellaSwag std: h0 vs h24 & 0.08 & --- & Negligible \\
\bottomrule
\end{tabular}
\end{table}

\section{Discussion}

\paragraph{The equalizer effect as retrieval-gated phenomenon.}
Our results position the equalizer effect as fundamentally tied to retrieval demand. Attention provides a direct structural advantage for retrieval tasks---the ability to look back over the full context to locate a target---that is relatively architecture-agnostic in its \emph{aggregate} benefit but architecture-specific in its \emph{per-example} mechanism. On tasks that primarily require commonsense reasoning or multi-step inference (HellaSwag, ARC-Challenge), attention's contribution is distributed across diverse reasoning pathways, yielding no systematic convergence pattern. The dose-response relationship across benchmarks ($\rho{=}{-}0.66$) is consistent with this account, though we acknowledge it is preliminary given $n{=}6$ benchmarks and $p{=}0.16$.

\paragraph{Scale as a necessary condition.}
The binary contrast between 340M and 1.3B demonstrates that sufficient model capacity is necessary for complete convergence. The 340M rebound on the hardest samples (62\% overlap with the bottom quintile) suggests capacity saturation: smaller models cannot simultaneously maintain recurrence quality and attention quality when attention layers dominate. However, the five-fold difference in training tokens (20B vs.\ 100B) prevents definitive attribution to parameter count alone. Scale is a necessary condition for complete convergence, likely through increased capacity to jointly optimize recurrence and attention, but the relative contributions of parameter count and training data remain an open question.

\paragraph{An opaque mechanism.}
Our systematic characterization---spanning surface features, baseline perplexity, topic distributions, and attention weight patterns---consistently finds that the architecture-specific nature of the equalizer effect is opaque to the features we examined. The GDN-specific perplexity signal ($p{=}0.031$, $d{=}0.19$) provides a suggestive lead but does not generalize across architectures and does not survive multiple comparison correction. Attention patterns show detectable but weak correlations ($r \approx 0.22$) that fall below our preregistered threshold. These negative results are informative: they rule out trivial explanations (``attention helps long-context samples'' or ``attention helps rare-word samples'') and establish that the mechanism is embedded in the recurrence-attention interaction at a computational level that standard and intermediate features do not capture. We characterize this as an architecture-specific effect whose mechanism is opaque---a negative result that constrains future mechanistic work.

\section{Limitations}

\paragraph{Architecture coverage.}
Our analysis spans three architectures (GDN, RetNet, HGRN2), representing three distinct recurrence mechanisms. We conducted an exhaustive search for additional hybrid architectures with comparable controlled training setups (Mamba hybrid, RWKV hybrid, Griffin) and found none with matching attention ratio sweeps and training conditions. While three architectures provide thousands of per-example data points, broader coverage would strengthen generalizability claims. We encourage the community to develop standardized hybrid model zoos spanning additional recurrence mechanisms.

\paragraph{Scale points.}
Only two scales (340M, 1.3B) are available from the model zoo. Our scale analysis is a binary contrast, not a multi-point trend, and we make no claims about extrapolation to larger scales.

\paragraph{Training data confound.}
The 340M and 1.3B models are trained on different amounts of data (20B vs.\ 100B tokens). While our analysis identifies capacity saturation as a plausible partial explanation for the 340M rebound, we cannot fully disentangle scale effects from training data effects.

\paragraph{Benchmark-level dose-response.}
The retrieval-demand gradient is estimated from $n{=}6$ benchmarks, yielding insufficient statistical power ($p{=}0.16$) for the cross-benchmark correlation. The dose-response relationship should be interpreted as a descriptive pattern, not an established quantitative law.

\section{Conclusion}

We present the first systematic quantification and decomposition of the equalizer effect in hybrid recurrent language models. Our per-example analysis reveals an equalizer paradox: attention drives aggregate performance convergence across architectures through almost entirely non-overlapping, architecture-specific mechanisms. This effect is retrieval-gated (centered on LAMBADA, with supporting evidence from ARC-Easy), scale-dependent (complete at 1.3B, partial at 340M), and mechanistically opaque to surface-level and intermediate features. These findings demonstrate that aggregate benchmark convergence can mask fundamentally different per-example behavior, suggesting that architecture selection decisions should be informed by the specific input distributions of target applications rather than aggregate leaderboard performance alone.

\end{document}
